\begin{abstract}
本道题目数据集为 50,000 条电动汽车消费者记录，三个子任务分别为对用户月度能源消耗量、月度充电成本与燃油替代支出的回归预测，评分指标为调整决定系数（adj-$R^2$）。

本题难点在于数据维度较高（23 列）、目标间存在级联关系，且题面提示的车型交互、复杂非线性等方向具有较强的误导性。在数据探索阶段，本文通过比率诊断与相关分析，识别出数据背后近似确定性的生成机制：能耗与周行驶里程成比例、燃油支出与日通勤里程成比例、充电成本近似等于能耗与电价的乘积。据此，本文并未采用复杂模型，而是以稀疏的物理公式与线性模型作为主方案。

在子任务一中，本文以周里程比率均值模型预测月度能耗，得到 adj-$R^2=0.876298$。在子任务二中，本文设计严格嵌套 OOF 级联方案，以任务一的折外预测乘以电价预测充电成本，得到 adj-$R^2=0.917817$，并给出真实能耗下的理论上界 0.9994。在子任务三中，本文以日通勤过原点 OLS 预测燃油支出，得到 adj-$R^2=0.905416$。

进一步的消融与残差诊断表明，题面建议的车型交互特征无增益，浅层 LightGBM 残差模型在 Bootstrap 置信区间下全面劣于稀疏基线，验证了``物理结构优先、复杂模型克制''的建模思路。本文的负结果报告为同类合成数据竞赛提供了可复用的方法学模板。

本文所有代码和中间文件均在支撑材料中随论文提交，并在 GitHub 仓库 \url{https://github.com/blackwhitetony/innoCupTaskB_Rematch} 中开源。核心实验均可在 CPU 上于数十秒内复现。

\textbf{关键词：}电动汽车 回归预测 稀疏物理模型 嵌套交叉验证 调整决定系数
\end{abstract}
